% ============================================================
%  TD Chapitre 5 : Mesures des angles et trigonometrie
% ============================================================
\clearpage
\chapter{Mesures des angles et trigonométrie}
\chaptermark{Mesures des angles et trigonométrie}

\begin{exercice}[conversions degrés-radians et longueur d'arc]
Convertir en radians les trois angles $210\dgr{}$, $240\dgr{}$ et $300\dgr{}$, en simplifiant chaque fraction, puis convertir en degrés l'angle $\dfrac{5\pi}{4}$ rad. Calculer ensuite, en centimètres, la longueur de l'arc découpé par un angle de $\dfrac{\pi}{4}$ rad sur un cercle de rayon $8$ cm. Contrôler chaque conversion avec la correspondance $180\dgr{} = \pi$ rad.
\end{exercice}

\begin{exercice}[Pythagore et trigonométrie]
\begin{enumerate}
\item une échelle de $10$ m est appuyée contre un mur, son pied étant à $4$ m du mur. À quelle hauteur le sommet de l'échelle touche-t-il le mur ?
\item un géomètre se place à $40$ m du pied d'une tour ; l'angle entre le sol horizontal et la ligne de visée du sommet vaut $25\dgr{}$. Calculer la hauteur de la tour, avec $\tan 25\dgr{} \approx 0{,}466$.
\end{enumerate}
\end{exercice}

\begin{exercice}[cercle trigonométrique, parité et périodicité]
\begin{enumerate}
\item donner les coordonnées $(\cos x;\ \sin x)$ du point image des angles $x = \dfrac{5\pi}{6}$ rad, en précisant le quadrant ;
\item en utilisant la parité, calculer $\cos(-\dfrac{\pi}{3})$ et $\sin(-\dfrac{\pi}{3})$ ;
\item en utilisant la périodicité, ramener $\sin\dfrac{7\pi}{3}$ à un angle du tableau des valeurs remarquables et donner sa valeur exacte.
\end{enumerate}
\end{exercice}

\begin{exercice}[équations trigonométriques]
Résoudre dans $\R$ les équations suivantes, en donnant l'ensemble des solutions avec le paramètre $k \in \Z$ :
\begin{enumerate}
\item $\cos x = \dfrac{\sqrt{2}}{2}$ ;
\item $\sin x = -\dfrac{1}{2}$ ;
\item $\cos x = 1$.
\end{enumerate}
\end{exercice}
